\section{Broader Impact}
\label{sec:impact}

The system we study is small (a 27-way style classifier on cached
features), but the design choices have broader implications. Three
points worth flagging.

First, our default ground-truth tree is built from a Western,
lineage-based art-history canon. Distance metrics computed against it
implicitly endorse that taxonomy. The Phase~3 ablation against
alternative trees is partial mitigation; deployment in a museum or
educational context should make the choice of tree explicit and
revisable rather than hard-coded.

Second, the WikiArt corpus is biased toward European painting and
sparsely represents non-Western traditions (Ukiyo-e is the only
non-Western top-level category). A retrieval system trained on this
data will under-rank non-Western works for ambiguous queries, and
hyperbolic geometry --- which we show preserves local hierarchy more
faithfully --- could amplify that effect by tightening neighborhoods
inside the dominant taxonomy. Users of any downstream system should
be informed of this bias.

Third, hyperbolic representations of style structure could in principle
be used to attribute or authenticate paintings. We do not endorse such
use without human-expert review: the classification accuracies reported
here ($\sim$50--65\,\% top-1 over 27 styles) are far below the
threshold required for legal or commercial decisions.
